\documentclass[12pt,a4paper]{article}
\usepackage[UTF8]{ctex}        % 中文支持
\usepackage{geometry}
\geometry{left=2.54cm,right=2.54cm,top=3.18cm,bottom=3.18cm}
\usepackage{enumerate}        % 列表格式
\usepackage{titlesec}         % 标题格式
\usepackage{xcolor}
\usepackage{indentfirst}
\setlength{\parindent}{2em}

%=========================导入宏包=========================
\usepackage{enumitem}



%=========================自定义功能区=========================

\newcommand{\textred}[1]{\textcolor{red}{#1}}  %红色字体
\newcommand{\textgreen}[1]{\textcolor{green}{#1}}  %绿色字体




\begin{document}



\begin{enumerate}
    
    \item   绝不可以胡思乱想；

    \item   没有发生的事情，不可在脑中演绎；

    \item   发生超过三天的事情，绝对不要去回味；

    \item   此时此刻，更值得你去感受和体会。

    \item   遇到问题时，脑中先想：“\textred{我该如何解决这个问题？}”
    
    \item   一件事情的成败不重要，重要的是你们曾经尝试过。\\[-3pt]
    {\small 要知道大多數人只是談論和夢想發財，而你們已經付出了行動。}

    \item   减少内耗：想太远未来、纠结过往都没用，只专注眼下这一件事。
    \item   降低期待，凡是尽力即可，不求完美，少和别人对比。
    \item   （睡前）闭眼清空杂念，只专注呼吸，脑子里冒想法别对抗，轻轻放过就好。
    \item   学习呢，其实就是，继续，不断地把自己的直觉再修正。如果可以修正到正确的方向，就可以比那些没有修正的人，多了一个思考的方法。






\end{enumerate}





\end{document}